% Diagrama da separação dos dados e da estimativa de desempenho (holdout, CV externa e CV interna).
% Requer as bibliotecas TikZ positioning e arrows.meta (as mesmas do pipeline).
\begin{figure}[!hbtp]
    \centering
    \resizebox{\columnwidth}{!}{%
    \begin{tikzpicture}[
        % Estilos (mesma família do pipeline)
        recntagle/.style={
            rectangle, rounded corners=3pt, inner sep=5pt, align=center,
            font=\footnotesize\bfseries, draw=black!70, fill=blue!10, line width=0.8pt
        },
        arrow/.style = {
            thick, ->, >={Stealth[length=2.5mm]}
        },
        box_delimitador/.style = {
            rectangle, rounded corners=8pt, draw=black!70, fill=black!2, line width=0.8pt
        },
        sub_block/.style = {
            rectangle, rounded corners=3pt, inner sep=4pt, align=center,
            draw=black!70, fill=blue!5, line width=0.7pt, font=\footnotesize, text width=5.5cm
        },
        result_block/.style = {
            sub_block, fill=green!15, draw=green!60!black, font=\footnotesize\bfseries
        },
        head_text/.style = {
            align=center, font=\footnotesize\bfseries
        },
        % Blocos das partições
        fold/.style = {
            rectangle, draw=black!70, line width=0.6pt, minimum height=0.8cm,
            inner sep=1pt, align=center, font=\scriptsize
        },
        ftrain/.style = {fold, fill=blue!18},
        fval/.style   = {fold, fill=green!25},
        ftest/.style  = {fold, fill=orange!35},
        fhold/.style  = {fold, fill=black!15}
    ]

        % ---------------- Painéis (desenhados primeiro, ficam ao fundo)
        \node [box_delimitador, minimum width=9.0cm, minimum height=15.15cm] at (-4.2,-10.775) {};
        \node [box_delimitador, minimum width=8.2cm, minimum height=15.15cm] at (4.8,-10.775) {};

        % ---------------- 1. Conjunto completo e divisão prévia
        \node (all) [recntagle, minimum width=7cm] at (0.3,0)
            {Conjunto completo de imagens\\[1pt]{\scriptsize\mdseries imagem RGB, rótulo e sequência topológica}};
        \node (dev) [recntagle, minimum width=9.0cm] at (-4.2,-2.4)
            {Desenvolvimento (80\%)\\[1pt]{\scriptsize\mdseries estimativa principal: validação cruzada aninhada}};
        \node (hold) [recntagle, fill=black!15, minimum width=8.2cm] at (4.8,-2.4)
            {\textit{Holdout} (20\%)\\[1pt]{\scriptsize\mdseries análise confirmatória, avaliada uma única vez}};

        \draw [arrow] (all.south) -- ++(0,-0.5) -| (dev.north);
        \draw [arrow] (all.south) -- ++(0,-0.5) -| (hold.north)
            node[pos=0.25, above, font=\scriptsize] {divisão estratificada (semente 42)};

        % ---------------- 2. Estimativa principal: CV aninhada
        \node (h_ext) [head_text, text width=8.4cm] at (-4.2,-3.9)
            {Validação cruzada externa: estratificada, 5 partições $\times$ 3 repetições = 15 partições};
        \draw [arrow] (dev.south) -- (h_ext.north);

        \foreach \x in {-7.25,-5.75,-4.25,-2.75} {
            \node [ftrain, minimum width=1.5cm] at (\x,-5.0) {Treino};
        }
        \node (t5) [ftest, minimum width=1.5cm] at (-1.25,-5.0) {Teste};
        % chave sob os blocos de treino
        \draw [thick, black!70] (-8.0,-5.5) -- (-8.0,-5.62) -- (-2.0,-5.62) -- (-2.0,-5.5);

        \node (h_int) [head_text, text width=5.8cm] at (-5.0,-6.9)
            {Validação cruzada interna: estratificada, 3 partições};
        \draw [arrow] (-5.0,-5.62) -- (h_int.north)
            node[midway, right, font=\scriptsize] {treino externo};

        \foreach \x in {-7.0,-5.0} {
            \node [ftrain, minimum width=2.0cm] at (\x,-7.9) {Treino interno};
        }
        \node [fval, minimum width=2.0cm] at (-3.0,-7.9) {Validação};

        \node (D) [sub_block, minimum height=3.1cm] at (-5.0,-10.45)
            {Para cada uma das 4 configurações ($\mathrm{lr}_{bb} \times \mathrm{lr}_{hd}$), treino com parada antecipada na validação interna. Escolhe-se a de menor entropia cruzada média, e o número de épocas é a média das épocas ótimas (mínimo 3)};
        \node (E) [sub_block, minimum height=1.7cm] at (-5.0,-13.35)
            {Ajuste final: treino em todo o treino externo, com número fixo de épocas e sem validação};
        \node (F) [sub_block, minimum height=1.3cm] at (-5.0,-15.2)
            {Avaliação única no teste externo: F1 macro, acurácia e AUC-ROC};
        \node (G) [result_block, minimum height=1.7cm] at (-5.0,-17.05)
            {Estimativa principal: média e desvio padrão do F1 macro entre as 15 partições};

        \draw [arrow] (-5.0,-8.3) -- (D.north);
        \draw [arrow] (D.south) -- (E.north);
        \draw [arrow] (E.south) -- (F.north);
        \draw [arrow] (F.south) -- (G.north);
        \draw [arrow] (t5.south) |- (F.east)
            node[pos=0.25, right, font=\scriptsize, align=left] {teste\\externo};

        % ---------------- 3. Análise confirmatória (holdout)
        \node (note) [rectangle, rounded corners=3pt, draw=black!70, dashed, fill=black!8, align=center,
                      text width=5.4cm, inner sep=4pt, font=\footnotesize] at (3.9,-4.7)
            {Isolado até a avaliação final: não participa da escolha de épocas, hiperparâmetros ou padronização};

        \node (h_all) [head_text, text width=5.8cm] at (3.9,-6.9) {Treino: desenvolvimento inteiro};
        \node (all80) [ftrain, minimum width=5.8cm] at (3.9,-7.9) {Desenvolvimento (80\% das imagens)};

        \node (R1) [sub_block, minimum height=3.1cm] at (3.9,-10.45)
            {Repete-se a seleção sobre todo o desenvolvimento: validação cruzada interna de 3 partições, as mesmas 4 configurações e número de épocas igual à média das épocas ótimas (mínimo 3)};
        \node (R2) [sub_block, minimum height=1.7cm] at (3.9,-13.35)
            {Ajuste final: treino em todo o desenvolvimento, com número fixo de épocas e sem validação};
        \node (R3) [sub_block, minimum height=1.3cm] at (3.9,-15.2)
            {Avaliação única no \textit{holdout} (20\%): F1 macro, acurácia e AUC-ROC};
        \node (R4) [result_block, minimum height=1.7cm] at (3.9,-17.05)
            {Resultado confirmatório: um valor por \textit{dataset}, sem medida de dispersão};

        \draw [arrow] (3.9,-8.3) -- (R1.north);
        \draw [arrow] (R1.south) -- (R2.north);
        \draw [arrow] (R2.south) -- (R3.north);
        \draw [arrow] (R3.south) -- (R4.north);
        \draw [arrow] (7.7,-2.95) -- (7.7,-15.2) -- (R3.east);
        \node [rotate=90, font=\scriptsize] at (8.05,-8.5) {\textit{holdout} (20\%)};

        % ---------------- Legenda
        \node [ftrain, minimum width=0.5cm, minimum height=0.3cm] at (-6.6,-19.15) {};
        \node [anchor=west, font=\scriptsize] at (-6.3,-19.15) {Treino};
        \node [fval, minimum width=0.5cm, minimum height=0.3cm] at (-4.4,-19.15) {};
        \node [anchor=west, font=\scriptsize] at (-4.1,-19.15) {Validação interna};
        \node [ftest, minimum width=0.5cm, minimum height=0.3cm] at (-0.9,-19.15) {};
        \node [anchor=west, font=\scriptsize] at (-0.6,-19.15) {Teste externo};
        \node [fhold, minimum width=0.5cm, minimum height=0.3cm] at (2.3,-19.15) {};
        \node [anchor=west, font=\scriptsize] at (2.6,-19.15) {\textit{Holdout}};

    \end{tikzpicture}%
    }
    \caption{Separação dos dados e estimativa do desempenho. Um conjunto de teste (\textit{holdout}) de 20\% é separado antes de qualquer treinamento. Os 80\% restantes (desenvolvimento) são avaliados por validação cruzada aninhada, com 15 partições externas (5 partições repetidas três vezes) e, em cada uma, uma validação cruzada interna de 3 partições para a escolha da configuração e do número de épocas. O mesmo procedimento, aplicado ao desenvolvimento inteiro, gera o modelo avaliado uma única vez no \textit{holdout}.}
    \label{fig:folds_tcc}
\end{figure}
